\chapter{The Logic of Financial Valuation}
\label{chap:The_Logic_of_Financial_Valuation}

\begin{openingcase}{The Same Expected Cash Flow, but the Same Value?}

An investor is comparing two one-year financial claims. Each has an expected payoff of RMB 100. Claim A pays RMB 120 when the economy is strong and RMB 80 when it is weak. Claim B does the opposite: it pays RMB 80 in the strong economy and RMB 120 in the weak economy. Suppose the two states are equally likely.

Their expected payoffs are identical:

\[
0.5(120)+0.5(80)=100.
\]

Yet the two claims need not have the same current value. Claim B pays more when economic conditions are weak, when an additional unit of cash may be particularly valuable. The distribution of payoffs across economic states can therefore affect valuation even when expected payoffs are identical.

\end{openingcase}

\begin{tcolorbox}[
  enhanced,
  colback=FinanceBlue!6,
  colframe=FinanceBlue,
  boxrule=0.6pt,
  sharp corners,
  left=8pt,
  right=8pt,
  top=7pt,
  bottom=7pt,
  before skip=1em,
  after skip=1.6em,
  title=\textbf{Core Questions},
  coltitle=white,
  colbacktitle=FinanceBlue,
  fonttitle=\bfseries,
  attach boxed title to top left={xshift=6pt,yshift=-2mm},
  boxed title style={
    boxrule=0pt,
    sharp corners,
    left=5pt,
    right=5pt,
    top=2pt,
    bottom=2pt
  }
]
\begin{enumerate}[
  leftmargin=*,
  label=\textcolor{FinanceBlue}{\textbf{\arabic*.}},
  itemsep=0.55em,
  topsep=0.2em
]

\item Why is expected cash flow alone insufficient to determine the value of a financial claim?

\item How do the timing, risk, and state dependence of future cash flows affect their present value?

\item How do information, contractual rights, and market prices interact in determining the value of financial claims?

\end{enumerate}
\end{tcolorbox}


\section{The Foundations of Financial Value}

A financial claim derives its value from the future economic benefits to which it entitles the holder. Before applying a valuation formula, three questions need to be addressed:


\begin{itemize}[
  leftmargin=2.2em,
  itemsep=0.45em,
  topsep=0.55em
]

\item What cash flows does the claim entitle its holder to receive?

\item When are these cash flows received?

\item How do these cash flows vary across future states of the world?

\end{itemize}

\subsection{Financial Assets as Claims on Future Economic Resources}

A financial asset represents rights to future economic resources, typically expressed through future cash flows.
A bond gives the holder contractual rights to specified payments, while common stock gives the holder a residual claim on the corporation after more senior claims have been satisfied. A derivative creates contingent cash flows that depend on an underlying asset, rate, index, or state variable.

The object of valuation is therefore the \emph{claim} defined by these rights, rather than simply the issuer, project, or underlying asset. Suppose a firm's operating assets generate cash flow $Y_t$ at date $t$. This cash flow may be allocated among several claimants:

\begin{equation}
Y_t=C_{D,t}+C_{E,t}+C_{O,t},
\label{eq:ch5-cash-flow-allocation}
\end{equation}
where $C_{D,t}$ denotes cash flow to debt holders, $C_{E,t}$ cash flow to equity holders, and $C_{O,t}$ cash flow to other contractual claimants.

The value of a financial claim depends not only on the cash flow generated by the underlying assets, but also on the contractual rules that determine how that cash flow is allocated.

Priority provides a simple illustration. Consider a firm financed only by debt and equity. Suppose it generates cash flow $Y$ in the next period and has debt with a promised repayment of $D$. Abstracting from other contractual provisions, debt holders receive

\begin{equation}
C_D=\min\{Y,D\},
\label{eq:ch5-debt-payoff}
\end{equation}
while equity holders receive

\begin{equation}
C_E=\max\{Y-D,0\}.
\label{eq:ch5-equity-payoff}
\end{equation}

Thus,

\begin{equation*}
C_D+C_E=Y.
\label{eq:ch5-claim-additivity}
\end{equation*}

The contractual structure does not change the firm's total cash flow; it determines how that cash flow is allocated across claimants for each realization of \(Y\). Figure~\ref{fig:ch5-debt-equity-payoff} illustrates this allocation.

\begin{figure}[H]

\centering

\begin{tikzpicture}[
=Latex,
axis/.style={-Latex, line width=0.8pt},
debt/.style={line width=1.0pt, FinanceBlue},
equity/.style={line width=1.0pt, dashed}
 ]

\draw[axis] (0,0) -- (7.2,0) node[right] {Firm cash flow $Y$};
\draw[axis] (0,0) -- (0,5.4) node[above] {Claim cash flow};

\draw[debt] (0,0) -- (3.6,3.6) -- (6.6,3.6);
\draw[equity] (0,0) -- (3.6,0) -- (6.6,3.0);

\draw[dashed, gray] (3.6,0) -- (3.6,4.6);

\node[below] at (3.6,0) {$D$};
\node[FinanceBlue, right] at (6.0,3.85) {Debt};
\node[right] at (6.2,2.55) {Equity};

\end{tikzpicture}
\caption{Allocation of firm cash flow between debt and equity}
\label{fig:ch5-debt-equity-payoff}

\end{figure}

\subsection{Time and the Present Value of Future Cash Flows}

A cash flow received in the future is not generally equivalent in value to an equal cash flow received today. Individuals may require compensation for postponing current consumption; funds available today can also be invested to generate future returns. These considerations provide the basic economic foundation for the time value of money and for interest and discounting.

A future cash flow can be expressed in terms of an equivalent amount today by discounting it at an appropriate rate.

Let \(V_0\) denote the present value of a claim that pays a certain cash flow \(CF_1\) one period from now. If the relevant one-period risk-free rate is \(r_{f}\), no-arbitrage equivalence requires:

\begin{equation*}
V_0(1+r_{f})=CF_1.
\label{eq:ch5-one-period-equivalence}
\end{equation*}
which gives

\begin{equation*}
V_0=\frac{CF_1}{1+r_{f}}.
\label{eq:ch5-one-period-pv}
\end{equation*}

The factor $1/(1+r_{f})$ is the one-period discount factor.
It gives the present value of one unit of cash flow
received one period in the future.

Assuming a constant per-period discount rate $r$,
the present value of a cash flow $CF_t$ received
$t$ periods from now is

\[
V_0=\frac{CF_t}{(1+r)^t}.
\]

For a sequence of future cash flows, the present
value is

\begin{equation}
V_0=\sum_{t=1}^{T}
\frac{CF_t}{(1+r)^t}.
\label{eq:ch5-pv-cash-flow-stream}
\end{equation}

Discounting expresses cash flows occurring at different dates in terms of value at a common date. Once expressed on this common basis, the cash flows can be compared and added.

\begin{figure}[H]
\centering
\begin{tikzpicture}[
  >=Latex,
  axis/.style={-Latex, line width=0.8pt},
  cash/.style={-Latex, line width=1.0pt, FinanceBlue},
  discount/.style={-Latex, dashed, line width=0.8pt}
]

\draw[axis] (0,0) -- (10.2,0) node[right]{Time};

\foreach \x/\t in {0/0,2/1,4/2,6/3,8/4}
{
  \draw (\x,0.12) -- (\x,-0.12);
  \node[below=4pt] at (\x,0) {$t=\t$};
}

\draw[cash] (2,0) -- (2,1.2) node[above] {$CF_1$};
\draw[cash] (4,0) -- (4,1.55) node[above] {$CF_2$};
\draw[cash] (6,0) -- (6,1.05) node[above] {$CF_3$};
\draw[cash] (8,0) -- (8,1.45) node[above] {$CF_4$};

\draw[discount] (2,2.25) -- (0.35,2.25);
\draw[discount] (4,2.25) -- (0.35,2.25);
\draw[discount] (6,2.25) -- (0.35,2.25);
\draw[discount] (8,2.25) -- (0.35,2.25);

\node[above] at (4.1,2.25) {Discount to $t=0$};

\end{tikzpicture}

\caption{Discounting future cash flows to the present}
\label{fig:ch5-discount-timeline}
\end{figure}

Several implications follow directly from discounting.

First, timing matters even when total nominal cash flow is unchanged. If two claims make the same total payments but one pays earlier, the earlier payments generally have greater present value when discount rates are positive.

Second, the sensitivity of present value to the discount rate increases with the time to payment. A change in the discount rate therefore has a larger effect on more distant cash flows.

Third, the timing of cash flows must be consistent with the discount-rate convention. An annual rate cannot be applied directly to quarterly cash flows without an appropriate conversion, and cash flows occurring at the beginning of a period must be distinguished from those occurring at the end.

\begin{example}{Timing Alone Changes Value}{timing-alone-changes-value}

Suppose Claim A pays RMB~1,000 one year from now and Claim B pays RMB~1,000 five years from now. If the relevant annual rate is 5 \%,

\[
V_A=\frac{1,000}{1.05}=952.38,
\]

while

\[
V_B=\frac{1,000}{1.05^5}=783.53.
\]

The nominal future payment is identical. The present value differs only because the payment dates differ.

\end{example}

\subsection{Risk and Required Return}

Time alone does not determine value. Many financial claims have risky future cash flows. Suppose there are two possible states next period. A claim pays $CF_H$ in one state and $CF_L$ in the other, with probabilities $p$ and $1-p$, respectively. Its expected cash flow is

\begin{equation}
\mathrm{E}[X]=pCF_H+(1-p)CF_L.
\label{eq:ch5-expected-cash-flow}
\end{equation}

Expected cash flow alone is generally insufficient for valuation. The value of a claim also depends on the risk of its cash flows. Cash flows that tend to be low when overall economic conditions are poor expose investors to greater risk. Such claims generally require a higher expected return. Two claims with the same expected payoff may consequently have different current values.

One common way to incorporate risk into valuation is through a required return. Let $r_f$ denote the risk-free rate and let $\pi$ denote the risk premium appropriate to the claim. The required return is

\begin{equation}
r=r_f+\pi.
\label{eq:ch5-required-return}
\end{equation}

A simplified one-period valuation based on expected cash flow is then

\begin{equation}
V_0=\frac{\mathrm{E}[CF]}{1+r}.
\label{eq:ch5-risk-adjusted-discount}
\end{equation}

The risk premium compensates investors for bearing risk that is relevant to valuation. Diversifiable risk may command little or no premium, while risk that tends to produce losses when overall economic conditions are poor may command a larger premium.

More generally, risk may be incorporated through risk-adjusted cash flows, discount rates, or state-dependent valuation. The treatment of risk should be consistent with the claim being valued.

\medskip

The main dimensions of financial valuation and their corresponding analytical inputs are summarized below.

\begin{table}[H]

\centering

\caption{Three questions in valuation}
\label{tab:ch5-three-valuation-questions}

\begin{tabularx}{\textwidth}{
  @{}
  >{\raggedright\arraybackslash\hyphenpenalty=10000\exhyphenpenalty=10000}p{3.3cm}
  Y
  Y
  @{}
}

\toprule

\textbf{Question}
& \textbf{Economic issue}
& \textbf{Typical valuation input} \\

\midrule

What is received?
& Rights to future economic resources
& Cash-flow definition and contractual payoff \\[20pt]

When is it received?
& Intertemporal valuation
& Discount factors or interest rates \\[20pt]

Under which states is it received?
& Risk and state dependence
& Risk premium or state-dependent valuation \\

\bottomrule

\end{tabularx}

\end{table}


\section{Valuation Frameworks}

The appropriate valuation method depends on the nature of the financial claim and the information available.

Discounted cash-flow (DCF) valuation derives value from expected cash flows attributable to a claim. Relative valuation uses observed prices of comparable assets, while contingent-claim valuation explicitly accounts for state-dependent payoffs, particularly when they are nonlinear or depend on future events.

Although these approaches differ in their inputs and implementation, each provides an estimate of value. Let $V$ denote the estimated value of a financial claim and $P$ its observed market price. In relative valuation, market prices of comparable claims serve as inputs for estimating the value of the claim under consideration. The resulting estimate can then be compared with its observed market price.

\subsection{Discounted Cash-Flow Valuation}

Discounted cash-flow valuation estimates the value of a claim by discounting the cash flows attributable to it.

If $r_j$ denotes the one-period required return applicable between dates $j-1$ and $j$, the value of a sequence of expected cash flows can be written as

\begin{equation}
V_0 = \sum_{t=1}^{T}
\frac{\mathrm{E}[CF_t]}
{\prod_{j=1}^{t}(1+r_j)}.
\label{eq:ch5-general-dcf}
\end{equation}

If the required return is constant at $r$,

\begin{equation}
V_0=\sum_{t=1}^{T}
\frac{\mathrm{E}[CF_t]}{(1+r)^t}.
\label{eq:ch5-constant-dcf}
\end{equation}

The formula is straightforward, but its application depends on how cash flows, risk, and long-run assumptions are defined.

\paragraph{Define the claim before forecasting the cash flow.}

Accounting earnings do not directly measure the cash flows available to a financial claim. Accrual accounting, noncash expenses, changes in working capital, and investment expenditure can cause reported earnings to differ from cash flows. The relevant cash-flow measure therefore depends on the claim being valued.

For a project, the relevant cash flow is the incremental
cash flow generated by undertaking the project. For
equity, it is the cash flow available to equity holders
after operating needs, investment, taxes, and payments
to more senior claimants. For debt, it consists of the
payments attributable to the debt claim, taking account
of the possibility of default and recovery.

\paragraph{Separate forecasting from valuation.}

Forecasting estimates future cash flows under specified assumptions; valuation translates those cash flows into present value. Keeping the two tasks distinct helps identify whether an error arises from the forecast itself or from the valuation of that forecast. A sophisticated discounting model cannot compensate for unrealistic growth assumptions, understated investment requirements, or inconsistencies in projected cash flows.

Suppose a project requires an initial investment $I_0$ and produces expected cash flows $\mathrm{E}[CF_t]$. Its net present value is

\begin{equation}
NPV=\sum_{t=1}^{T}
\frac{\mathrm{E}[CF_t]}{(1+r)^t}-I_0.
\label{eq:ch5-npv}
\end{equation}

A positive NPV means that, under the stated assumptions, the present value of expected future cash flows exceeds the resources required today. It indicates positive economic surplus when those cash flows are valued at the relevant required return; it does not imply that losses cannot occur in particular realized states.

\begin{example}{From Cash Flows to NPV}{from-cash-flows-to-npv}


A project requires an initial investment of RMB~10 million and is expected to generate RMB~3.0 million at the end of year 1, RMB~4.0 million at the end of year 2, and RMB~5.5 million at the end of year 3. Suppose the required return is 8\%.

The present value of the future cash flows is

$$
PV=\frac{3.0}{1.08}+\frac{4.0}{1.08^2}+\frac{5.5}{1.08^3}=10.57
$$
million RMB, approximately. Therefore,

$$
NPV=10.57-10.00=0.57
$$

million RMB.

The project has a positive NPV under the stated cash-flow and required-return assumptions.

\end{example}

\paragraph{Terminal value extends the forecast beyond the explicit horizon.}

Many assets continue to generate cash flows after the explicit forecast period. Discounted cash-flow valuation therefore often separates value into two components: cash flows forecast individually through date $T$, and a terminal value at date $T$ representing the value of cash flows expected after that date.

\begin{equation}
V_0=\sum_{t=1}^{T}
\frac{\mathrm{E}[CF_t]}{(1+r)^t}+\frac{TV_T}{(1+r)^T}.
\label{eq:ch5-dcf-terminal}
\end{equation}

If cash flow after year $T$ is expected to grow perpetually at a constant rate $g$, the terminal value at date $T$ is

\begin{equation}
TV_T=\frac{CF_{T+1}}{r-g},
\qquad r>g.
\label{eq:ch5-growing-perpetuity}
\end{equation}

Terminal value is sensitive to both $r$ and $g$. When the explicit forecast period is short, terminal value may constitute a substantial share of total estimated value. The valuation then depends heavily on the long-run assumptions underlying these parameters.

The condition $r>g$ is required for the growing-perpetuity formula to yield a finite value. The assumed long-run growth rate should also be consistent with the sustainable growth of the relevant market or economy.

\paragraph{Sensitivity analysis identifies the assumptions that matter most.}

A valuation is conditional on its assumptions. Sensitivity analysis examines how estimated value changes in response to changes in key inputs.

Sensitivity analysis measures the dependence of valuation results on selected assumptions rather than the likelihood of alternative outcomes. More distant cash flows are particularly sensitive to the discount rate, while terminal values are often sensitive to long-run growth assumptions.

The inputs to a discounted cash-flow valuation are based on economic, accounting, contractual, and market analysis. Forecasts of sales, margins, reinvestment, and other cash flows reflect assumptions about future operating conditions, while required returns depend on the risk of the claim and relevant market evidence. DCF combines these inputs in a consistent valuation framework.

\subsection{Relative Valuation: Comparables and Market Multiples}

Relative valuation estimates the value of an asset from the observed prices of economically comparable assets.

Suppose comparable firms trade at price-to-earnings ratios around 15. If the firm being valued has earnings per share of RMB~2, the implied price is $15\times 2=30$.

A valuation multiple expresses market value relative to earnings, book value, sales, cash flow, or another economic measure. Common examples include

\begin{equation*}
\frac{P}{E},
\qquad
\frac{P}{B},
\qquad
\frac{EV}{EBITDA},
\qquad
\frac{EV}{Sales}.
\label{eq:ch5-common-multiples}
\end{equation*}

The numerator and denominator of a valuation multiple should refer to the same set of claimants. Market value of equity $P$ is matched with measures attributable to equity holders, such as earnings or book equity. Enterprise value $EV$ is matched with operating measures available to all capital providers, such as EBITDA or sales.

Valuation multiples vary across firms with differences in their underlying economic characteristics. Consider a share whose dividends are expected to grow at a constant rate $g$ indefinitely. If the next-period dividend per share is $D_1$ and the required return on equity is $r$, its value is

\begin{equation}
P_0=\frac{D_1}{r-g}.
\label{eq:ch5-dividend-growth}
\end{equation}

This valuation relation is known as the Gordon growth model.
With a constant dividend payout ratio $\alpha$, next-period dividends per share satisfy $D_1=\alpha E_1$, where $E_1$ denotes next-period earnings per share. Substituting this relation into Equation~\eqref{eq:ch5-dividend-growth} gives:

\begin{equation}
\frac{P_0}{E_1}=\frac{\alpha}{r-g}.
\label{eq:ch5-pe-drivers}
\end{equation}

Even in this simplified setting, the price--earnings multiple depends on payout, required return, and expected growth. More generally, growth depends on both reinvestment and the return earned on that reinvestment, while profitability affects the cash flows generated from the firm's asset base. Valuation multiples therefore reflect underlying assumptions about growth, risk, payout, and profitability.

The reliability of relative valuation depends on the economic comparability of the assets and the measurement of the variables used in the multiple. Industry classification provides a useful starting point, but firms within the same industry may differ in growth, leverage, operating risk, accounting practices, and geographic exposure. Relative valuation incorporates current market prices and can provide a useful benchmark for assessing estimated values, but market-wide mispricing, accounting differences, cyclical or negative earnings, and differences in financial leverage can reduce the usefulness of observed multiples.


\subsection{Contingent-Claim Valuation: State-Dependent Payoffs}

Some financial claims have payoffs that vary across future states. Their valuation therefore depends not only on expected cash flows but also on how these payoffs are distributed across states.

Options provide a clear example, but state-dependent payoffs arise more broadly. Equity reflects limited liability, credit instruments are exposed to default states, and insurance contracts depend on specified loss events. Similar features appear in guarantees, convertible securities, and structured claims.

Suppose there are $S$ possible states at date $T$, and let $CF_T(s)$ denote the claim's payoff at date $T$ if state $s$ occurs. Let $q_s$ denote the price at date 0 of a security that pays one unit of currency at date $T$ if state $s$ occurs and zero otherwise. The price $q_s$ is called the state price for state $s$.

State prices reflect the value of receiving resources in different states. A unit of payoff received in a state in which resources are relatively scarce may be more valuable than the same unit received in a more favorable state. State-dependent valuation therefore assigns different values to payoffs according to the states in which they occur, rather than treating all units of expected payoff as economically equivalent.

Under state-price valuation, the value of the claim at date 0 is

\begin{equation}
V_0=\sum_{s=1}^{S}q_s CF_T(s).
\end{equation}

\begin{example}{State-Dependent Valuation}{state-dependent-valuation}

Consider two equally likely states, $s=1,2$. Claim A pays 120 in state 1 and 80 in state 2, while Claim B pays 80 in state 1 and 120 in state 2. Both claims have an expected payoff of 100.

Suppose the corresponding state prices are

$$
q_{1}=0.40,
\qquad
q_{2}=0.55.
$$

The values of the two claims at date 0 are

\begin{align}
V_A
&=
0.40(120)+0.55(80)=92,\\
V_B
&=0.40(80)+0.55(120)=98.
\end{align}

Claim B has the higher value at date 0 because it pays more in state 2, where the state price is higher.

\end{example}

State prices are also related to risk-free discounting. A claim that pays one unit in every state is risk free. Its  value at date 0 is

\begin{equation}
\sum_{s=1}^{S}q_s=\frac{1}{1+r_f},
\label{eq:ch5-state-price-riskfree}
\end{equation}
where $r_f$ is the one-period risk-free rate. Thus, the sum of state prices equals the one-period risk-free discount factor. More generally, state-price valuation expresses the value of a state-dependent claim at date 0 as the sum of its payoffs weighted by their corresponding state prices.

State-price valuation is closely related to the principle of no arbitrage. Let $CF_T^A(s)$ and $CF_T^B(s)$ denote the payoffs of portfolios A and B, respectively, in state $s$ at date $T$, and let $P_A$ and $P_B$ denote their prices at date 0. Under the no-arbitrage principle, identical payoffs in every possible state require identical prices at date 0:

\begin{equation}
CF_T^A(s)=CF_T^B(s)
\quad \text{for every } s
\quad \Longrightarrow \quad
P_A=P_B.
\label{eq:ch5-law-one-price}
\end{equation}

If their prices differed, an investor could buy the lower-priced portfolio and sell the higher-priced portfolio while receiving offsetting future payoffs, thereby creating an arbitrage opportunity.

State prices reflect both the time value of money and the relative value of payoffs across different states. A claim is valued by weighting each state-contingent payoff by its corresponding state price. 

The state-price representation also clarifies the nature of risk in valuation. Payoffs in states with higher state prices contribute more to value at date 0 than equal payoffs in states with lower state prices. Consequently, payoff variance alone does not determine value. Even for claims with the same expected payoff, greater payoff variability does not necessarily imply a lower value; the effect depends on how payoffs are distributed across states with different state prices.

\section{Information, Contracts, and Value Formation}

Valuation is influenced by the information available about future cash flows and risk, as well as by the contractual rights attached to the claim being valued.

\subsection{Information and Cash-Flow Expectations}

Valuation depends on expectations about future cash flows and risk, and these expectations are formed from available information. Revenue, default losses, investment needs, and long-run growth are not known in advance; they are estimated from information available at the valuation date.

Let $\mathcal{I}_0$ denote the information available at date 0. The value of a claim can be written as

\begin{equation}
V_0=\sum_{t=1}^{T}
\frac{\mathrm{E}[CF_t\mid \mathcal{I}_0]}
{\prod_{j=1}^{t}(1+r_j)},
\label{eq:ch5-information-conditional-value}
\end{equation}
where expected cash flows are conditional on the available information and $r_j$ denotes the required return for period $j$.

New information can change value by altering expected
cash flows, required returns, or both. Let
$\mathcal{I}_0'$ denote an updated information set
available at date 0. Then,

\begin{equation*}
\mathrm{E}[CF_t\mid \mathcal{I}_0']
\neq
\mathrm{E}[CF_t\mid \mathcal{I}_0]
\end{equation*}

New information affects valuation through two channels: it can change expected future cash flows or alter the perceived risk of the claim and, consequently, its required return.

The effect of additional information on value is not
unidirectional. It may increase or decrease expected
cash flows and alter assessments of underlying risk.
Improved information quality can reduce estimation
uncertainty, although it does not necessarily reduce
the underlying economic risk. Valuation depends on
the relevance and reliability of information as well
as on its quantity.

\subsection{Contracts and the Allocation of Cash-Flow Rights}

Contracts define the cash-flow rights associated with a financial claim, including the conditions governing payment and the priority of the claim relative to others. These contractual features are directly relevant to valuation.

For a firm financed only by debt and equity, let $V_A$ denote the value of its assets. The values of the corresponding claims satisfy

\begin{equation}
V_A=V_D+V_E,
\label{eq:ch5-value-allocation}
\end{equation}
where $V_D$ and $V_E$ denote the values of debt and equity, respectively. Contractual provisions can affect how this value is distributed between the claims even when the value of the underlying operating assets is unchanged.

Seniority determines payment priority, collateral affects recovery in default, and guarantees redistribute losses across parties.

\begin{example}{Collateral and Claim Value}{collateral-claim-value}


Suppose two one-year debt claims each promise repayment of 100. The first is secured by collateral with a realizable value of 60 in default, while the second is unsecured. If default occurs and the firm's remaining assets are insufficient to repay both claims in full, the secured claim has a higher recovery. The two claims can therefore have different values even though their promised payments and maturities are identical.

\end{example}

Contracts also influence value through incentives. Covenants may limit additional borrowing, dividend payments, or other actions that increase the risk borne by creditors. Such restrictions can protect creditor claims, although they may reduce operating flexibility.

Other contractual provisions create state-contingent rights. A convertible bond, for example, gives the holder the right to exchange the bond for equity under specified conditions. Embedded options such as conversion or call provisions alter the payoffs that apply in different future states.

Contractual provisions affect claim value by determining payment priority, influencing incentives, and defining state-dependent payoffs. For complex financial instruments, valuation therefore depends on the contractual rights attached to the claim rather than on the instrument's label alone.

\subsection{Market Prices and Valuation}

\paragraph{Market prices and price discovery.}
A market price is the price at which a claim trades, whereas valuation estimates the claim's value based on specified information, assumptions, and pricing relationships. Market price and estimated value are related but conceptually distinct.

In liquid markets, prices reflect the interaction of participants' information, beliefs, and trading orders, through which new information may be incorporated into prices. Market prices do not, however, provide an assumption-free measure of fundamental value. Investors may differ in information, beliefs, constraints, tax positions, funding costs, and risk exposures. Prices may also be affected by liquidity needs, leverage, market-making capacity, forced transactions, and temporary order imbalances.

\paragraph{Market prices as a valuation benchmark.}
Market prices provide an external benchmark for valuation. A substantial difference between an estimated value and the observed market price may reflect differences in cash-flow expectations, assessments of risk, discount-rate assumptions, omitted factors, or market mispricing. Such a difference does not by itself establish that either the valuation or the market price is incorrect; it indicates that the underlying assumptions and sources of disagreement require further examination.

No-arbitrage imposes a stronger restriction on relative prices. If two traded assets have identical payoffs in every possible state at the same future date and relevant market frictions are absent, they must have the same price at date 0. Otherwise, an arbitrage opportunity would exist. This principle, known as the law of one price, links the prices of securities with equivalent payoffs and is fundamental to fixed-income and derivative valuation.

\paragraph{Valuation when replication is limited.}
For claims that cannot be closely replicated by traded securities, valuation relies more heavily on cash-flow forecasts and required-return assumptions. Private firms and long-horizon investment projects, for example, may lack close market equivalents. Prices of comparable firms and traded claims, together with interest rates, credit spreads, and option prices, nevertheless provide useful valuation inputs.

\begin{figure}[H]
\centering

\begin{tikzpicture}[
  >=Latex,
  box/.style={
    draw,
    rounded corners=2pt,
    align=center,
    minimum height=0.9cm,
    inner xsep=12pt,
    inner ysep=5pt
  },
  flow/.style={
    -Latex,
    line width=0.8pt,
    FinanceBlue
  }
]

% Top nodes
\node[box] (info) at (-2.0,2.2)
{Information and\\expectations};

\node[box] (contract) at (2.0,2.2)
{Contractual\\cash-flow rights};

% Central sequence
\node[box] (valuation) at (0,0.3)
{Valuation assessments};

\node[box] (trading) at (0,-1.6)
{Trading and price discovery};

\node[box] (price) at (0,-3.5)
{Observed market price};

% Main flows: 90-degree turns
\coordinate (val-left) at
  ($(valuation.north west)!0.09!(valuation.north east)$);
\draw[flow]
  (info.south) -- (val-left);

\coordinate (val-right) at
  ($(valuation.north west)!0.91!(valuation.north east)$);

\draw[flow]
  (contract.south) -- (val-right);


% Vertical flows
\draw[flow]
  (valuation.south) -- (trading.north);

\draw[flow]
  (trading.south) -- (price.north);

% Feedback from observed market price to valuation
\draw[flow]
  (price.west)
  -- ++(-2.0,0)
  -- node[left, midway, align=center, xshift=-3pt]
     {market\\evidence}
     ($(valuation.west)+(-2.0,0)$)
  -- (valuation.west);

\end{tikzpicture}

\caption{Information, contractual rights, valuation, and market prices}
\label{fig:ch5-value-formation}

\end{figure}

Figure~\ref{fig:ch5-value-formation} summarizes these interactions. Information affects expectations about future cash flows and risk, while contracts define the cash-flow rights attached to each claim. These factors enter valuation, which in turn informs trading decisions. Trading and price discovery generate observed market prices, which provide market evidence for assessing and revising valuations.

Figure~\ref{fig:ch5-value-formation} summarizes these interactions. Information shapes expectations about future cash flows and risk, while contractual terms define the cash-flow rights attached to each claim. Valuation combines these inputs and informs trading decisions. Trading and price discovery produce market prices, which in turn provide evidence for assessing and revising valuations.

\begin{tcolorbox}[
  enhanced,
  breakable,
  colback=white,
  colframe=FinanceBlue,
  boxrule=0.7pt,
  sharp corners,
  borderline west={2.2pt}{0pt}{FinanceBlue},
  title={Financial Perspective: Valuation as a Conditional Statement},
  coltitle=white,
  colbacktitle=FinanceBlue,
  fonttitle=\bfseries,
  before skip=1.2em,
  after skip=1.2em
]

A valuation estimate is conditional on the contractual cash-flow rights, the information available, the timing assumptions, and the treatment of risk used in the analysis. Changes in these inputs may change the estimated value. A sound valuation states its assumptions explicitly and applies them consistently, allowing the result to be evaluated and revised as information changes.

\end{tcolorbox}

%========================================================
% Summary
%========================================================

\TailModule{summary}{Summary}

Financial valuation translates future cash flows into current value. The object being valued is a financial claim, so valuation begins by identifying the cash flows and contractual rights attached to that claim. Payment priority, collateral, conversion rights, and other contractual provisions can cause claims on the same underlying assets to have different values.

The timing of cash flows matters, as cash flows occurring at different dates must be converted to values at a common date. Risk also affects value: two claims with the same expected payoff can have different values when their payoffs are distributed differently across future states. Required returns provide one way to incorporate risk into valuation, while state-price valuation makes the dependence of value on future states explicit.

Discounted cash-flow valuation combines forecasts of claim-specific cash flows with appropriate discount rates. Its reliability depends on consistent treatment of cash flows and required returns, realistic forecasts, and economically plausible terminal assumptions. Relative valuation uses the observed prices of comparable assets, but valuation multiples reflect differences in growth, risk, profitability, payout, and other fundamentals. Contingent-claim valuation values state-contingent payoffs using state prices and no-arbitrage relationships.

Information, contracts, and valuation are closely connected. Information affects expectations about future cash flows and risk, while contracts define the rights attached to each claim. Trading and price discovery generate observable market prices, which provide evidence for evaluating and revising valuations. A valuation is conditional on the information, contractual rights, timing assumptions, and treatment of risk used in the analysis. Its reliability depends on internal consistency, economic plausibility, and its relationship to available market evidence.

%========================================================
% Key Terms
%========================================================

\TailModule{keyterms}{Key Terms}

Financial valuation; financial claim; cash-flow right; present value; discount factor; required return; risk-free rate; risk premium; expected cash flow; discounted cash-flow valuation; net present value; terminal value; growing perpetuity; sensitivity analysis; relative valuation; comparable firm; market multiple; price--earnings ratio; enterprise value; contingent claim; state-contingent payoff; state price; law of one price; no-arbitrage; price discovery; market price

%========================================================
% Financial AI Lab
%========================================================

\TailModule{aiprompt}{Financial AI Lab}

\textbf{Lab: Auditing an AI-Generated Valuation}

Choose a publicly traded company and ask an AI system to produce a simplified discounted cash-flow valuation using publicly available information. Require the AI to state the claim being valued, the forecast horizon, the cash-flow measure, the growth assumptions, the terminal-value method, and the discount rate.

Audit the valuation in five stages.

\begin{enumerate}

\item \textbf{Identify the claim.}
Determine whether the valuation refers to equity or to the entire enterprise, and whether the cash flows and discount rate are defined consistently with that claim.

\item \textbf{Verify the calculations.}
Reproduce the present-value calculations independently. Check units, cash-flow timing, compounding, and the treatment of terminal value.

\item \textbf{Evaluate the economic assumptions.}
Examine whether the cash-flow forecasts are economically plausible, whether long-run growth can reasonably persist, whether reinvestment is consistent with growth, and whether the discount rate is appropriate for the risk of the claim.

\item \textbf{Conduct sensitivity analysis.}
Identify two key assumptions to which the estimated value is particularly sensitive and report how changes in those assumptions affect the valuation.

\item \textbf{Compare the valuation with market evidence.}
Compare the result with the observed market price and one simple relative-valuation measure. Ensure that the quantities being compared refer to the same claim. Do not treat either the AI estimate or the market price as necessarily correct; explain the assumptions or market factors that may account for the difference.

\end{enumerate}

The final submission should identify at least one issue in the AI-generated valuation that is not apparent from the final valuation estimate alone. The issue may involve calculation, economic assumptions, internal consistency, or verification against market evidence.

%========================================================
% Applied Problems
%========================================================

\TailModule{appprobs}{Applied Problems}

\begin{enumerate}

\item \textbf{Timing and present value.}

A claim pays RMB~5,000 in three years. The relevant annual discount rate is 6 \%.

\begin{enumerate}

\item Calculate the present value of the claim.

\item Recalculate the present value if the payment occurs in six years instead.

\item Explain why postponing the payment reduces its present value when the discount rate is positive.

\end{enumerate}


\item \textbf{NPV and required return.}

A project requires an initial investment of RMB~12 million and is expected to produce cash flows of RMB~3 million, RMB~4 million, RMB~5 million, and RMB~4 million at the end of years 1 through 4.

\begin{enumerate}

\item Calculate the project's NPV at a required return of 7 \%.

\item Recalculate the NPV at a required return of 11 \%.

\item Explain why the project's value changes even though its expected cash flows remain unchanged.

\end{enumerate}


\item \textbf{Terminal-value sensitivity.}

A business is expected to generate cash flow of RMB 20 million in year 6, the first year after the explicit forecast period. From year 7 onward, cash flow is expected to grow at a constant annual rate \(g\) indefinitely. Assume that the year-6 cash flow remains unchanged across all scenarios.

\begin{enumerate}

\item Calculate the terminal value at the end of year 5 if $r=9\%$ and $g=3\%$.

\item Recalculate the terminal value if $g=4\%$, holding $r$ unchanged.

\item Recalculate the terminal value if $r=8\%$ and $g=3\%$.

\item Compare the results in parts (b) and (c). Explain why a one-percentage-point increase in $g$  and a one-percentage-point decrease in $r$ produce the same terminal value under the stated assumptions.

\end{enumerate}


\item \textbf{Relative valuation and consistency.}

Firm A has an enterprise value of RMB~8 billion and EBITDA of RMB~800 million. Firm B has EBITDA of RMB~500 million.

\begin{enumerate}

\item Calculate Firm A's EV/EBITDA multiple.

\item Apply this multiple to Firm B to estimate its implied enterprise value.

\item Suppose Firm B has RMB~1.2 billion of net debt. Calculate its implied equity value.

\item Explain why an enterprise-value multiple should not be applied directly to an earnings measure attributable only to equity holders.

\end{enumerate}


\item \textbf{State-price valuation.}

There are two equally likely future states. Claim A pays 140 in the strong state and 60 in the weak state. Claim B pays 80 in the strong state and 120 in the weak state.

\begin{enumerate}

\item Calculate the expected payoff of each claim.

\item Suppose the state prices are $q_S=0.38$ and $q_W=0.56$. Calculate the  values of each claim at date 0.

\item Compare the expected payoffs and the values of two claim at date 0.

\item Explain why equal expected payoffs need not imply equal values at date 0.

\end{enumerate}


\item \textbf{Contractual allocation of value.}

A firm will generate cash flow of either 150 or 70 next period, with each state equally likely. It has debt promising repayment of 100.

\begin{enumerate}

\item Calculate the payoffs to debt and equity in each state.

\item Verify that the debt and equity payoffs sum to the firm's total cash flow in each state.

\item Using state prices $q_H=0.40$ and $q_L=0.55$, calculate the current values of debt, equity, and the firm as a whole.

\item Verify that the values of debt and equity sum to the value of the firm.

\end{enumerate}

\end{enumerate}

%========================================================
% Further Reading
%========================================================

\TailModule{furtherread}{Further Reading}

\begin{itemize}

\item Fisher, Irving (1930), 
\emph{The Theory of Interest: As Determined by Impatience to Spend Income and Opportunity to Invest It}, 
Macmillan.

\item Williams, John Burr (1938), 
\emph{The Theory of Investment Value}, 
Harvard University Press.

\item Modigliani, Franco, and Merton H. Miller (1958), 
``The Cost of Capital, Corporation Finance and the Theory of Investment,'' 
\emph{The American Economic Review}, 48(3), 261--297.

\item Sharpe, William F. (1964), 
``Capital Asset Prices: A Theory of Market Equilibrium under Conditions of Risk,'' 
\emph{The Journal of Finance}, 19(3), 425--442.

\item Damodaran, Aswath (2012), 
\emph{Investment Valuation: Tools and Techniques for Determining the Value of Any Asset}, 
3rd ed., Wiley.

\end{itemize}

\ChapterTail{summary,keyterms,aiprompt,appprobs,furtherread}
